\section{\texttt{av-robustbench}}
\label{sec:bench}

To make every reported number reproducible and to give the community a common interface for AV robust detection, we release \texttt{av-robustbench}, a Python library that organizes the evaluation logic of this paper. It provides: (i) a detector-adapter interface so any frozen-encoder detector can be plugged in; (ii) attacks --- PGD-$\ell_\infty$ and PGD-$\ell_2$ with joint AV projection, plus Square-Attack and AutoAttack-style wrappers; (iii) randomized-smoothing certification (isotropic and PCA-aligned anisotropic); (iv) a degradation battery (social re-encoding, H.264, JPEG, audio corruptions); and (v) JSON ``robustness cards'' and a leaderboard schema. Following the RobustBench philosophy~\cite{croce2021robustbench} but adapted to multimodal forensic threat models, its role here is reproducibility: every attack, certificate, degradation, and transfer result in this paper maps to a structured entry rather than an ad-hoc script, so the full study can be regenerated end to end.

\section{Limitations}
\label{sec:limitations}

We highlight five concrete limitations with future directions. (1) \emph{Feature-space scope}: CertAV certifies the frozen representation, not raw pixels or waveforms; composing the certificate with an encoder Lipschitz bound to obtain an end-to-end input guarantee is open. (2) \emph{Single primary dataset}: results center on FakeAVCeleb with transfer to LAV-DF; broader evaluation across manipulation families and datasets is needed, and the $61.4\%$ transfer certified accuracy shows headroom. (3) \emph{Class imbalance}: the $1{:}10$ real/fake ratio depresses conditional coverage on the real class (Table~\ref{tab:conformal}); class-conditional conformal calibration is a natural fix. (4) \emph{Computational cost}: each certified clip needs $n=1000$ forward passes, standard for smoothing but limiting for real-time use; caching frozen features mitigates but does not remove this. (5) \emph{Subspace-invariance approximation}: Theorem~\ref{thm:scaling} assumes exact invariance, whereas $\cos^2\theta\approx0.18\neq0$ indicates small residual off-manifold sensitivity; tightening the theory to bounded (rather than exact) invariance would make the amplification factor a guarantee rather than an idealized upper bound.

\section{Conclusion}
\label{sec:conclusion}

CertAV makes certified robustness practical for audio-visual deepfake detection by placing the guarantee at the level of frozen multimodal features. Its central finding is geometric and predictive: the intrinsic dimension of the representation --- not the training procedure --- governs certifiability (Theorem~\ref{thm:scaling}), which explains the large isotropic radii, predicts encoder behavior (Corollary~\ref{cor:amp}), and motivates manifold-aware anisotropic smoothing for a $3.4\times$ on-manifold gain with zero abstention. Robust conformal prediction completes the picture with a coverage guarantee for every sample. Together these results turn a striking empirical observation into a coherent, certified story, and establish representation geometry as a first-class design axis for trustworthy multimodal forensics.
